%% FullCourse_10Lecs -- Lecture 9, Part 9 Appendix
%% Assembled by tools/lec09_assemble.py from reorg_manifest.yaml: Phase A of
%% Lec09_Reorg_Plan_2026-09-12.md as amended by Lec09_Reorg_Plan_Review_2026-09-12.md.
%% Each "% ---- <ID> · <TAG> · TIER" marker names a frame's source deck and line;
%% keep the markers until Phase B is done (lec09_assemble.py check reads them).
%% Owns: the interoperability stack (MCP and A2A explained only here), frameworks, reasoning models, scheduled runtimes, computer use
%% Owns: dated install notes per harness, keybindings and templates, dated comparisons (incl. model-specific task horizons), resources, Context Dilemma references
%% Defers: the harness file mapping -> T2b (review A2.5)

%% Shared preamble for the FullCourse_10Lecs topic decks.
%%
%% Each topic .tex \input's this file, then defines its own \title, \date
%% override (if any), \graphicspath, and \begin{document}.
%%
%% Reconstructed verbatim from the inlined copy preserved in
%% Lec10_Case_Studies/Lec10_T8_Case6_DeepRamsey.tex (lines 23-190), which is
%% explicitly annotated there as "from ../shared_preamble.tex, copied verbatim".
%% Base preamble matches MiniCourse_8hr/Slides/shared_preamble.tex; this version
%% additionally provides \sectiondivider, the conceptbox/cautionbox/examplebox/
%% takeawaybox tcolorboxes, and \compacttables used across the split decks.

\documentclass[10pt,english,aspectratio=169]{beamer}

\usepackage{ae,aecompl}
\usepackage[T1]{fontenc}
\usepackage{textcomp}
\usepackage[utf8]{inputenc}
\usepackage{babel}
\usepackage{datetime2}

\setcounter{secnumdepth}{3}
\setcounter{tocdepth}{3}

\usepackage{amsmath, amssymb, amsthm, graphicx}
\usepackage{booktabs, multirow}
\usepackage{tikz}
\usepackage{pgfplots}
\pgfplotsset{compat=1.16}
\usepackage[most]{tcolorbox}
\tcbuselibrary{skins, breakable, theorems}
\usepackage{listings}
\usepackage{xcolor}

\lstset{
  basicstyle=\ttfamily\small,
  keywordstyle=\color{blue!80!black}\bfseries,
  commentstyle=\color{green!50!black}\itshape,
  stringstyle=\color{orange!80!black},
  backgroundcolor=\color{gray!8},
  frame=single,
  rulecolor=\color{gray!40},
  breaklines=true,
  showstringspaces=false,
  numbers=none,
}

\lstdefinestyle{pythonstyle}{
    language=Python,
    basicstyle=\ttfamily\footnotesize,
    keywordstyle=\color{blue!80!black}\bfseries,
    commentstyle=\color{green!50!black}\itshape,
    stringstyle=\color{orange!80!black},
    frame=single,
    breaklines=true,
    showstringspaces=false,
    numbers=left,
    numbersep=5pt,
    numberstyle=\tiny\color{gray},
    backgroundcolor=\color{gray!8},
    rulecolor=\color{gray!40}
}

\ifx\hypersetup\undefined
  \AtBeginDocument{%
    \hypersetup{bookmarks=false,breaklinks=true,pdfborder={0 0 0},colorlinks=true,
      linkcolor=DarkText,urlcolor=RedTitle}%
  }
\else
  \hypersetup{bookmarks=false,breaklinks=true,pdfborder={0 0 0},colorlinks=true,
    linkcolor=DarkText,urlcolor=RedTitle}%
\fi

\makeatletter
\newcommand\makebeamertitle{%
  \begin{frame}[plain]
    \vspace*{1.0cm}
    \centering
    {\usebeamerfont{title}\usebeamercolor[fg]{title}\inserttitle\par}
    \vspace{1.0cm}
    {\usebeamerfont{author}\usebeamercolor[fg]{author}\insertauthor\par}
    \vspace{0.2cm}
    {\usebeamerfont{date}\usebeamercolor[fg]{date}\insertdate\par}
  \end{frame}}%
\AtBeginDocument{%
  \let\origtableofcontents=\tableofcontents
  \def\tableofcontents{\@ifnextchar[{\origtableofcontents}{\gobbletableofcontents}}
  \def\gobbletableofcontents#1{\origtableofcontents}%
}

\usetheme{metropolis}
\usefonttheme{professionalfonts}

\definecolor{LightGray}{RGB}{230,230,230}
\definecolor{RedTitle}{RGB}{0,0,200}
\definecolor{VeryLightGray}{RGB}{245,245,245}
\definecolor{DarkText}{RGB}{0,0,0}

\setbeamercolor{background canvas}{bg=white}
\setbeamercolor{title}{fg=RedTitle}
\setbeamercolor{author}{fg=DarkText}
\setbeamercolor{date}{fg=DarkText}
\setbeamercolor{frametitle}{bg=VeryLightGray,fg=DarkText}
\setbeamercolor{normal text}{fg=DarkText}
\setbeamerfont{title}{size=\large}

\setbeamersize{text margin left=5mm, text margin right=5mm}
\addtobeamertemplate{frametitle}{}{\vspace{-0.5em}}
\newcommand{\reducevspace}{\vspace{-0.5cm}}

% Shared section divider and box styles for split topic decks.
\newcommand{\sectiondivider}[2]{%
\begin{frame}[plain,noframenumbering]
\begin{center}
\vfill
{\Large\textbf{#1}\par}
\vspace{0.35cm}
{\normalsize #2\par}
\vfill
\end{center}
\end{frame}
}

\newtcolorbox{conceptbox}[1][]{
  colback=blue!3,
  colframe=RedTitle!55,
  boxrule=0.35mm,
  arc=1mm,
  left=2mm,
  right=2mm,
  top=1mm,
  bottom=1mm,
  #1
}
\newtcolorbox{cautionbox}[1][]{
  colback=red!3,
  colframe=red!50!black,
  boxrule=0.35mm,
  arc=1mm,
  left=2mm,
  right=2mm,
  top=1mm,
  bottom=1mm,
  #1
}
\newtcolorbox{examplebox}[1][]{
  colback=green!3,
  colframe=green!45!black,
  boxrule=0.35mm,
  arc=1mm,
  left=2mm,
  right=2mm,
  top=1mm,
  bottom=1mm,
  #1
}
\newtcolorbox{takeawaybox}[1][]{
  colback=VeryLightGray,
  colframe=RedTitle!55,
  boxrule=0.35mm,
  arc=1mm,
  left=2mm,
  right=2mm,
  top=1mm,
  bottom=1mm,
  #1
}
\newcommand{\compacttables}{\renewcommand{\arraystretch}{1.15}}

\setbeamertemplate{navigation symbols}{}
\setbeamertemplate{footline}{%
  \hfill%
  \usebeamercolor[fg]{page number in head/foot}%
  \usebeamerfont{page number in head/foot}%
  \insertframenumber\,/\,\inserttotalframenumber\kern1em\vskip2pt%
}

\makeatother

\author{Zhigang Feng}
% "date with time": compile timestamp so each build is identifiable.
\date{\today\ \DTMcurrenttime}


% Suppress redundant auto-generated metropolis section page; \sectiondivider frames are the dividers we keep.
\metroset{sectionpage=none}

\graphicspath{{./pic/}{./figures/}{../}}

\usetikzlibrary{positioning,arrows.meta,shapes,calc,fit,backgrounds}


\title[AI for Econ Research]{AI for Economic Research: Dynamic Models, Language, and Agents\\
Lecture 9 Appendix: The Landscape (Dated), Setup Details, References}

\begin{document}

\makebeamertitle

% ---- T1b-F01 · EDIT · TIER: READ · from T1b:24
% TODO(Phase B): Replace the agenda list with the appendix contents; keep the [solid]/[hype] reading convention.
\begin{frame}{What This Appendix Holds}
\textbf{Everything in this appendix is dated: re-check it before teaching. The main decks point here instead of repeating product facts.}

\vspace{0.2cm}

\begin{enumerate}
  \item \textbf{The interoperability stack} --- MCP, A2A, AGNTCY/ACP, neutral governance
  \item \textbf{Orchestration patterns} --- frameworks, context engineering, the multi-agent cost-benefit
  \item \textbf{Reasoning models \& inference-time compute} --- think budgets, measured time horizons
  \item \textbf{Autonomous coding \& computer-use agents} --- scheduled runtimes, sandboxing
  \item \textbf{Benchmarks \& measurement validity} --- ties to the course's measurement theme
  \item \textbf{Safety \& reproducibility} --- untrusted inputs, chain-of-custody
\end{enumerate}

\vspace{0.15cm}
\footnotesize\textit{Reading convention: \textbf{[solid]} $=$ primary or citable; \textbf{[hype]} $=$ vendor self-report (counter-source given).}

\vspace{0.05cm}
{\scriptsize \textit{Facts current as of June 2026. Model versions are stated with release dates; the field moves monthly --- recheck before teaching.}}
\end{frame}


%=============================================================================
\section{The Landscape (Dated)}
%===========================================================

\sectiondivider{The Landscape (Dated)}{Compete on models, cooperate on plumbing --- as of September 2026}

% ---- T1b-F02 · KEEP · TIER: READ · from T1b:51
\begin{frame}{Three Layers of Agent Interoperability}
\textbf{The 2025--26 story is standard-setting: vendors compete on models but increasingly cooperate on the wiring between agents and tools.}

\vspace{0.18cm}

\begin{center}
\footnotesize
\renewcommand{\arraystretch}{1.2}
\begin{tabular}{|p{2.6cm}|p{3.0cm}|p{4.9cm}|}
\hline
\textbf{Layer} & \textbf{Protocol} & \textbf{What it standardizes} \\ \hline
Agent $\leftrightarrow$ tool & \textbf{MCP} (Model Context Protocol) & how an agent discovers and calls external tools/data sources \\ \hline
Agent $\leftrightarrow$ agent & \textbf{A2A} (Agent-to-Agent) & how independent agents delegate tasks to one another \\ \hline
Discovery / identity & \textbf{AGNTCY / ACP} & agent directories, capability cards, identity and trust \\ \hline
\end{tabular}
\end{center}

\vspace{0.18cm}

\begin{itemize}\footnotesize
    \item MCP was introduced by Anthropic (Nov 2024); it is now \emph{multi-vendor}, not proprietary.
    \item \textbf{OpenAI adopted MCP} at DevDay (Oct 6, 2025); Google and Microsoft tooling support it too. \textbf{[solid]}
    \item Governance moved to a neutral home: the \textbf{Linux Foundation Agentic AI Foundation (AAIF)}, formed Dec 9, 2025, now hosts MCP and related projects. \textbf{[solid]}
\end{itemize}
\end{frame}

% ---- T1b-F03 · KEEP · TIER: READ · from T1b:77
\begin{frame}{Why Neutral Governance Is a Network-Effects Story}
\textbf{For economists, the interoperability stack is a textbook standard-setting / network-effects case.}

\vspace{0.22cm}

\begin{itemize}
    \item A shared protocol lowers switching costs: a tool wrapped once in MCP works across Claude, ChatGPT, Gemini, and open agents
    \item Network effects make the standard more valuable as adoption grows --- which is why a single-vendor protocol faces adoption resistance
    \item Moving MCP under the \textbf{Linux Foundation AAIF} (Dec 9, 2025) is the classic move: hand the ``plumbing'' to a neutral body so rivals will build on it
    \item Competition then concentrates where rents remain: model quality, latency, and price --- not the connectors
\end{itemize}

\vspace{0.22cm}

\begin{takeawaybox}
\footnotesize\textbf{Reproducibility note.} Protocols version. Pin the \emph{spec version} you used (e.g.\ MCP spec date, A2A version) alongside the model version and date --- a tool call that worked under one spec revision may behave differently under the next.
\end{takeawaybox}
\end{frame}

% ---- T1b-F05 · KEEP · TIER: READ · from T1b:125
\begin{frame}{Frameworks Reached 1.0 (2025--26)}
\textbf{The orchestration layer matured from research demos to versioned, supported frameworks.}

\vspace{0.16cm}

\begin{center}
\footnotesize
\renewcommand{\arraystretch}{1.15}
\begin{tabular}{|p{4.6cm}|p{3.0cm}|p{2.6cm}|}
\hline
\textbf{Framework} & \textbf{Milestone} & \textbf{Vendor} \\ \hline
LangGraph 1.0 & Oct 2025 & LangChain \\ \hline
OpenAI Agents SDK ($+$ AgentKit) & SDK Mar 2025 & OpenAI \\ \hline
Google Agent Dev. Kit (ADK) & v1.0 May 2025 & Google \\ \hline
Microsoft Agent Framework & 1.0 Apr 2026 & Microsoft \\ \hline
CrewAI & 1.0 & CrewAI \\ \hline
Anthropic \textbf{Claude Agent SDK} & Sep 2025 & Anthropic \\ \hline
\end{tabular}
\end{center}

\vspace{0.12cm}
\footnotesize\textbf{Naming note.} What was the \emph{Claude Code SDK} was renamed the \textbf{Claude Agent SDK} on Sep 29, 2025, reflecting generalization beyond coding. \textbf{[solid]}

\vspace{0.05cm}
\textbf{Durable lesson:} the patterns outlive the framework. Learn orchestrator--worker once; map it onto whichever SDK your institution adopts.
\end{frame}

% ---- T1b-F08 · KEEP · TIER: READ · from T1b:199
\begin{frame}{Reasoning Models and Think-Budget Control}
\textbf{Frontier models now spend variable inference-time compute, and increasingly call tools \emph{inside} the reasoning loop.}

\vspace{0.16cm}

\begin{itemize}\footnotesize
    \item \textbf{Tool use in the loop.} OpenAI o3 / o4-mini (Apr 2025) interleave tool calls with chain-of-thought rather than reasoning first and acting after.
    \item \textbf{Think-budget routers.} GPT-5 (Aug 2025) $\to$ GPT-5.5 (Apr 2026) route between fast and deep modes per query.
    \item \textbf{Toggleable extended thinking + configurable token budget.} \textbf{Claude Opus 4.5} (Nov 2025) lets you turn extended thinking on/off and cap the thinking-token budget; it was the first model past \textbf{80\% on SWE-bench Verified}, and Opus pricing was cut to \$5/\$25 per Mtok (input/output).
    \item \textbf{Parallel deep reasoning.} Gemini ``Deep Think'' explores multiple reasoning paths in parallel.
\end{itemize}

\vspace{0.12cm}
\begin{cautionbox}
\footnotesize\textbf{Benchmark caveat (carry this everywhere).} The 80\% SWE-bench Verified figure is real but the metric is contaminated/saturating: OpenAI \emph{stopped reporting} SWE-bench Verified in Feb 2026. Pair any Verified headline with SWE-bench Pro ($\sim$55--59\% at the frontier, held-out). \textbf{[solid]}
\end{cautionbox}
\end{frame}

% ---- T1b-F10 · KEEP · TIER: READ · from T1b:248
\begin{frame}{Coding Agents Became Scheduled, Autonomous Runtimes}
\textbf{In 2025--26, coding agents shifted from interactive assistants to background runtimes with persistent state.}

\vspace{0.16cm}

\begin{itemize}\footnotesize
    \item \textbf{OpenAI Codex ``Goal Mode''} became the default (May 2026), with persistent cross-session state --- the agent resumes a goal rather than starting fresh each chat.
    \item \textbf{Anthropic ``Code with Claude'' (May 6, 2026)} reframed Claude Code as an \emph{orchestrator}:
    \begin{itemize}\scriptsize
        \item \textbf{Auto Mode} with prompt-injection / destructive-action classifiers $+$ human approval gates
        \item \textbf{Routines} triggered by cron or webhook
        \item \textbf{Worktrees} for parallel branches
        \item sandboxed \textbf{Managed Agents} with scoped credentials
    \end{itemize}
\end{itemize}

\vspace{0.12cm}
\begin{takeawaybox}
\footnotesize\textbf{Lesson for research infra.} Before an agent runs unattended on real data or repos, configure sandboxing, credential scoping, and approval gates first. ``Autonomous'' is a permissions decision, not a model capability.
\end{takeawaybox}
\end{frame}

% ---- T1b-F11 · KEEP · TIER: READ · from T1b:270
\begin{frame}{Computer-Use Agents: Strong Benchmarks, Volatile Products}
\textbf{Browser/computer-use agents got good on paper but the standalone products churned.}

\vspace{0.16cm}

\begin{itemize}\footnotesize
    \item Computer-use / browser agents reached $\sim$\textbf{78--84\% on OSWorld-Verified} (human baseline $\sim$72\%).
    \item \textbf{But standalone products are volatile:} Google \emph{killed} Project Mariner (May 2026); OpenAI folded Operator into ChatGPT Agent / Atlas.
\end{itemize}

\vspace{0.12cm}
\begin{cautionbox}
\footnotesize\textbf{Flag the benchmark numbers.} OSWorld scores are \textbf{scaffold-dependent and largely self-reported}: the same model scores very differently under different harnesses, and ``above human baseline'' does not mean reliable on \emph{your} unseen workflow. Treat as directional. \textbf{[solid claim that products churned; benchmark figures = vendor/scaffold-dependent]}
\end{cautionbox}

\vspace{0.12cm}
\textbf{Lesson:} do not build durable research infrastructure on a single experimental product. Build on the protocol layer (MCP) and your own files; treat any one browser agent as replaceable.
\end{frame}

% ---- T1b-F16 · EDIT · TIER: READ · from T1b:388
\begin{frame}{Synthesis: The Landscape Reinforces the Workflow Spine}
\textbf{Every module above maps back onto the terminal-agent workflow of T1--T5d and onto the course's measurement / identification theme.}

\vspace{0.16cm}

\begin{center}
\footnotesize
\renewcommand{\arraystretch}{1.18}
\begin{tabular}{|p{3.2cm}|p{7.2cm}|}
\hline
\textbf{Landscape module} & \textbf{Connection to the spine} \\ \hline
Interoperability (MCP/A2A) & build on the neutral protocol layer, not one product; pin spec versions \\ \hline
Orchestration & T5b's sub-agent review \emph{is} orchestrator--worker; pay for multi-agent only when value $>$ token cost \\ \hline
Reasoning \& time horizons & ``runs for hours'' $\to$ a measured, reliability-thresholded number \\ \hline
Coding / computer-use & autonomy is a sandboxing + approval-gate decision, set before the run \\ \hline
Benchmarks & a headline score is not a measurement until you know its validity (Verified vs.\ Pro) \\ \hline
Safety \& reproducibility & untrusted-input discipline $+$ logged chain-of-custody (dates, versions, seeds) \\ \hline
\end{tabular}
\end{center}

\vspace{0.1cm}
\begin{center}
\textbf{The plumbing changes monthly; the research discipline does not.}\\
\footnotesize\textit{Compete on questions and judgment; cooperate on standards; measure before you trust.}
\end{center}
\end{frame}


%=============================================================================
\section{Setup Details}
%===========================================================

\sectiondivider{Setup Details}{Dated install notes, keys, and templates}

% ---- T6-F13 · EDIT · TIER: READ · from T6:275
% TODO(Phase B): Review A2: dated install notes for each tested harness (Claude Code, Gemini CLI, Codex CLI), from CLI_FACTS_2026-09.md.
\begin{frame}[fragile]{Appendix A: Installing One Terminal Agent (Dated)}
\textbf{Three operational steps, illustrated with Claude Code on macOS:}

\vspace{0.1cm}
\begin{enumerate}
    \item \textbf{Install the runtime} (Node.js for Claude Code; Python for Aider; analogous for others):
\begin{lstlisting}[style=pythonstyle,language={},basicstyle=\ttfamily\tiny,numbers=none]
brew install node  # or download from nodejs.org (v18+)
\end{lstlisting}
    \item \textbf{Install the agent CLI:}
\begin{lstlisting}[style=pythonstyle,language={},basicstyle=\ttfamily\tiny,numbers=none]
npm install -g @anthropic-ai/claude-code   # Claude Code
# pip install codex-cli                    # OpenAI Codex CLI
# pip install aider-chat                   # Aider
\end{lstlisting}
    \item \textbf{Verify:} \texttt{claude --version}, then \texttt{cd your-project-folder \&\& claude}.
\end{enumerate}

\vspace{0.1cm}
\footnotesize\textbf{Mac + OneDrive note:} always \texttt{cd} into the project folder before launching the agent. Running from a \texttt{/Volumes/\ldots} path can trigger permission errors.

\vspace{0.05cm}
\textbf{Dated note:} subscription tiers, model IDs, and install commands change every few months. Recheck the official docs before teaching or running the lab.
\end{frame}

% ---- T6-F14 · EDIT · TIER: READ · from T6:300
% TODO(Phase B): Review A2.6: point to the lab's Step 0 setup check, which routes each student by harness.
\begin{frame}{Appendix A2: Setup for the Lab}
\textbf{Lightweight setup checklist for a terminal-agent workflow}

\vspace{0.2cm}

\begin{itemize}
    \item Install one terminal agent and confirm it launches
    \item Create a project folder with a project-specific system prompt (e.g., \texttt{CLAUDE.md}, \texttt{AGENTS.md}, \texttt{GEMINI.md}), \texttt{scripts/}, and \texttt{output/}
    \item Make sure Python or R runs from the terminal
    \item Keep the first session small: one file, one check, one output
\end{itemize}

\vspace{0.25cm}

\textbf{Dated note:} exact setup steps, plans, and rate limits change quickly. Recheck the official docs before the lab.
\end{frame}

% ---- T6-F15 · KEEP · TIER: READ · from T6:317
\begin{frame}[fragile]{Appendix B: Keybindings and Minimal Templates}
\begin{columns}[T]
\begin{column}{0.48\textwidth}
\textbf{Useful commands}

\vspace{0.15cm}
\footnotesize
\begin{tabular}{|l|l|}
\hline
Compact session & \texttt{/compact} \\ \hline
Clear session & \texttt{/clear} \\ \hline
Interrupt & \texttt{Esc} \\ \hline
Slash menu & \texttt{/} \\ \hline
\end{tabular}
\end{column}
\begin{column}{0.48\textwidth}
\textbf{Minimal project-system-prompt}
\begin{lstlisting}[style=pythonstyle,language={},basicstyle=\ttfamily\tiny,numbers=none]
# CLAUDE.md / AGENTS.md / GEMINI.md
Project: ...
Goal: ...
Language: ...
Checks: ...
\end{lstlisting}
\end{column}
\end{columns}
\end{frame}

% ---- T6-F16 · EDIT · TIER: READ · from T6:345
% TODO(Phase B): Re-verify every entry; add the model-specific task horizons moved out of T1 (review A5).
\begin{frame}{Appendix C: Dated Comparisons and Resources}
\textbf{As of September 2026, examples of terminal-agent workflows include:}
\begin{itemize}
    \item Claude Code
    \item Codex CLI
    \item Gemini CLI
    \item Aider
\end{itemize}

\vspace{0.25cm}

\textbf{Good places to recheck before teaching:}
\begin{itemize}
    \item official Claude Code documentation
    \item official OpenAI Codex CLI documentation
    \item course notes and lab instructions
\end{itemize}

\vspace{0.2cm}

\textbf{Teaching rule:} keep dated product comparisons in the appendix, not in the conceptual core of the lecture.
\end{frame}


%=============================================================================
\section{Resources}
%===========================================================

\sectiondivider{Resources}{Docs, repositories, and references}

% ---- NEW · TIER: READ
% TODO(Phase B): Official docs for each tested harness; Sant'Anna and Blattman skill repositories; Goldsmith-Pinkham's guides; METR; OWASP agentic Top-10; CORE-Bench; the lab.
\begin{frame}{Resources}
\textbf{Placeholder --- written in Phase B.} The specification is in the source comment above this frame.
\end{frame}

% ---- T6-F17 · KEEP · TIER: READ · from T6:368
\begin{frame}{Appendix D: Context Dilemma --- References}
\scriptsize

\textbf{Context as implicit prior}
\begin{itemize}
    \item Xie, Raghunathan, Liang, Ma (2022). An Explanation of In-context Learning as Implicit Bayesian Inference. ICLR. arXiv:2111.02080.
\end{itemize}

\textbf{Sycophancy and RLHF limits}
\begin{itemize}
    \item Sharma et al.\ (2023). Towards Understanding Sycophancy in Language Models. arXiv:2310.13548.
    \item Perez et al.\ (2023). Discovering Language Model Behaviors with Model-Written Evaluations. ACL Findings. arXiv:2212.09251.
    \item Casper et al.\ (2023). Open Problems and Fundamental Limitations of RLHF. TMLR. arXiv:2307.15217.
    \item Wolf et al.\ (2023). Fundamental Limitations of Alignment in Large Language Models. arXiv:2304.11082.
\end{itemize}

\textbf{Self-correction and adversarial review}
\begin{itemize}
    \item Huang et al.\ (2024). Large Language Models Cannot Self-Correct Reasoning Yet. ICLR. arXiv:2310.01798.
    \item Valmeekam et al.\ (2023). Can Large Language Models Really Improve by Self-critiquing Their Own Plans? NeurIPS. arXiv:2310.08118.
    \item Du et al.\ (2023). Improving Factuality and Reasoning through Multiagent Debate. arXiv:2305.14325.
    \item Irving, Christiano, Amodei (2018). AI Safety via Debate. arXiv:1805.00899.
\end{itemize}

\textbf{Monoculture and the economist workflow}
\begin{itemize}
    \item Kleinberg \& Raghavan (2021). Algorithmic Monoculture and Social Welfare. PNAS 118(22).
    \item Bommasani et al.\ (2022). Picking on the Same Person: Does Algorithmic Monoculture Lead to Outcome Homogenization? NeurIPS. arXiv:2211.13972.
    \item Korinek (2023). Language Models and Cognitive Automation for Economic Research. NBER WP 30957; JEL.
\end{itemize}
\end{frame}

%===========================================================
\end{document}
